\subsection*{Question 2.1 - Force and torque as functions of the rotor thrusts}

The goal is to write the total external force $f$ and torque $n$ acting on the quadrotor, both expressed in the body frame $\{B\}$, as linear functions of the four rotor thrusts $f_T$ and of the gravity force $f_g$.

\paragraph{Conventions.}
We use the frame and numbering of Fig.~1 of the handout.
The axes $x_B$, $y_B$ and $z_B$ form a right-handed frame with $z_B$ pointing down, so every rotor thrust $f_i$ points along $-z_B$.
Let $r > 0$ denote the distance from the center of mass to each rotor (the arm length).
Reading the position of the rotors off the figure, the rotors lie along the body axes at
\begin{equation*}
r_1 = \begin{bmatrix} r \\ 0 \\ 0 \end{bmatrix}, \qquad
r_2 = \begin{bmatrix} 0 \\ -r \\ 0 \end{bmatrix}, \qquad
r_3 = \begin{bmatrix} -r \\ 0 \\ 0 \end{bmatrix}, \qquad
r_4 = \begin{bmatrix} 0 \\ r \\ 0 \end{bmatrix}.
\end{equation*}

\paragraph{Step 1 - Force.}
Two things push on the vehicle: gravity and the four rotors.
Gravity is naturally given in the inertial frame, $f_g = m g\, e_3$.
Since $R(\lambda)$ rotates vectors from $\{B\}$ to $\{I\}$, its transpose expresses an inertial vector in $\{B\}$, so the gravity force in $\{B\}$ is $R^T(\lambda) f_g$.
Each rotor produces a force $-f_i e_3$ that is already written in $\{B\}$.
Adding the four contributions gives the total thrust force
\begin{equation*}
-\sum_{i=1}^{4} f_i\, e_3 = -e_3 \begin{bmatrix} 1 & 1 & 1 & 1 \end{bmatrix} f_T .
\end{equation*}
Therefore
\begin{equation*}
f = \underbrace{R^T(\lambda)}_{M} f_g + \underbrace{\Big(-e_3 \begin{bmatrix} 1 & 1 & 1 & 1 \end{bmatrix}\Big)}_{N} f_T ,
\end{equation*}
that is,
\begin{equation*}
M = R^T(\lambda), \qquad
N = \begin{bmatrix} 0 & 0 & 0 & 0 \\ 0 & 0 & 0 & 0 \\ -1 & -1 & -1 & -1 \end{bmatrix}.
\end{equation*}
The matrix $M$ depends on the attitude because gravity is defined in $\{I\}$ and its direction, seen from the body, changes as the vehicle rotates.
The matrix $N$ is constant because the thrusts always act along the same body axis.

\paragraph{Step 2 - Torque from the thrusts.}
A force $F_i$ applied at the position $r_i$ produces the torque $r_i \times F_i$ about the center of mass.
With $F_i = (0, 0, -f_i)^T$ and $r_i = (r_{x}, r_{y}, 0)^T$ we have
\begin{equation*}
r_i \times F_i = \begin{bmatrix} -f_i\, r_{y} \\ f_i\, r_{x} \\ 0 \end{bmatrix}.
\end{equation*}
Evaluating this for each rotor:
\begin{align*}
\text{rotor 1: } & (0,\ r f_1,\ 0), &
\text{rotor 2: } & (r f_2,\ 0,\ 0), \\
\text{rotor 3: } & (0,\ -r f_3,\ 0), &
\text{rotor 4: } & (-r f_4,\ 0,\ 0).
\end{align*}
Summing the contributions, the thrusts produce a roll torque $n_x = r(f_2 - f_4)$ and a pitch torque $n_y = r(f_1 - f_3)$, and no yaw torque.

\paragraph{Step 3 - Torque from the rotor drag.}
Besides the thrust, every spinning rotor exerts a reaction torque $n_i = c f_i$ on the airframe about $z_B$.
The sign of this torque depends on the spin direction.
In Fig.~1 rotors 1 and 3 spin one way and rotors 2 and 4 the other way (the counter-rotating pairs), which gives
\begin{equation*}
n_z = c\,(f_1 - f_2 + f_3 - f_4).
\end{equation*}
Opposite spin directions are exactly what allows the vehicle to produce a yaw torque by speeding up one pair and slowing down the other.

\paragraph{Step 4 - Assembling $P$.}
Collecting the three components,
\begin{equation*}
n = \begin{bmatrix} n_x \\ n_y \\ n_z \end{bmatrix} = P f_T, \qquad
P = \begin{bmatrix} 0 & r & 0 & -r \\ r & 0 & -r & 0 \\ c & -c & c & -c \end{bmatrix}.
\end{equation*}
In summary, equations (2) and (3) hold with
\begin{equation*}
M = R^T(\lambda), \qquad
N = -e_3 \begin{bmatrix} 1 & 1 & 1 & 1 \end{bmatrix}, \qquad
P \text{ as above.}
\end{equation*}
